% ===========================================================================
%  12 — Catalogo de experiencias
% ===========================================================================
\chapter{Catálogo de experiências}
\label{cap:experiencias}

O catálogo alvo é de 48 experiências. Cinco existem, e são elas que definiram o
gabarito que as próximas seguem. Este capítulo documenta cada uma até onde a
mecânica impõe exigência ao \emph{software} --- e não como material pedagógico.

\section{Visão do catálogo}

\begin{tabularx}{\linewidth}{@{}L{28mm}L{14mm}L{12mm}L{20mm}X@{}}
\toprule
\headrow \thd{Experiência} & \thd{Séries} & \thd{Público} & \thd{Capabilities} & \thd{Módulo na API}\\
\midrule
\cd{exp-01-instanooria} & \cd{8EF} & aluno &
  telemetria, progresso, ia &
  \arq{motor}, \arq{geracao}, \arq{atividades}, \arq{perfis}, \arq{dossie}\\
\zebra \cd{exp-02-unootube} & \cd{4EF} & aluno &
  telemetria, progresso & \arq{unootube}\\
\cd{exp-70-vila-do-lume} & \textbf{transversal} & aluno &
  telemetria & \arq{vila}\\
\zebra \cd{exp-90-plano-de-aula} & --- & \textbf{professor} &
  telemetria, ia & \arq{planos}\\
\cd{exp-91-bem-estar} & --- & \textbf{professor} &
  telemetria, ia & \arq{bemestar}\\
\zebra \arq{\_template} & --- & --- & telemetria, progresso & ---\\
\bottomrule
\end{tabularx}

\begin{tabularx}{\linewidth}{@{}L{28mm}L{18mm}X@{}}
\toprule
\headrow \thd{Experiência} & \thd{Porta em \emph{dev}} & \thd{Onde aparece no hub}\\
\midrule
\cd{exp-01-instanooria} & \cd{5173} & \rota{/exp/exp-01-instanooria}\\
\zebra \cd{exp-02-unootube} & \cd{5174} & \rota{/exp/exp-02-unootube}\\
\cd{exp-90-plano-de-aula} & \cd{5175} & \rota{/professor/plano-de-aula}\\
\zebra \cd{exp-91-bem-estar} & \cd{5176} & \rota{/professor/bem-estar}\\
\cd{exp-70-vila-do-lume} & \cd{5177} & \rota{/exp/exp-70-vila-do-lume}\\
\bottomrule
\end{tabularx}

\begin{nota}
As portas vivem em \arq{apps/hub/lib/experiencias-dev.ts} --- uma por
experiência. É esse mapa que o hub usa quando não há manifesto, e é por isso que
\env{EXP\_INTERNAL\_HOST} tem de \emph{preservar a porta}
(Seção~\ref{sec:manifesto}).
\end{nota}

\section{exp-01 --- Instanooria}

Clone de rede social de imagens. O aluno navega um \emph{feed} curado por um
motor adaptativo que o empurra progressivamente para um dos dois lados de uma
controvérsia; na virada da missão, o dossiê mostra o que ele não viu.

\begin{tabularx}{\linewidth}{@{}L{36mm}X@{}}
\toprule
\headrow \thd{Exigência ao \emph{software}} & \thd{Como é atendida}\\
\midrule
\emph{Feed} adaptativo por blocos, determinístico & \arq{motor/curadoria.ts}:
  função pura que recebe um gerador de aleatoriedade. É isso que permite ao teste
  afirmar \enquote{só engaja A \implica{} domínio $>$ 0,75 até o bloco 4}.\\
\zebra Serialização por sessão & \cd{pg\_advisory\_xact\_lock(runId)} dentro da
  transação, com \cd{timeout} de \SI{20}{\second}.\\
Gabarito de curadoria fora do cliente & \arq{motor/sanitizador.ts} remove
  \cd{lado}, \cd{perfil} e \cd{intensidade} antes de o \emph{feed} sair.\\
\zebra Geração de conteúdo em lote, assíncrona & \arq{geracao/}: fila em
  \cd{exp\_01\_generation\_jobs}, \emph{worker} \emph{in-process}, reciclagem de
  órfão por \cd{updated\_at}.\\
Momento coletivo controlado pelo professor & \cd{activities.etapa} é um relógio
  \textbf{coletivo}: o professor abre e fecha, e a tela do aluno trava
  conforme.\\
\zebra Reação ao vivo, desligável sem \emph{deploy} & \cd{activities.ia\_ao\_vivo},
  alterável por \rota{PATCH /atividades/:id/ia}.\\
Largada por perfil, sorteio ou redistribuição & \arq{perfis/largada.ts},
  com precedência explícita e \cd{lado\_origem} gravado para auditoria.\\
\zebra Casca de aparelho de 375$\times$812 & \cd{viewport} no manifesto \implica{}
  \cd{DispositivoFrame}.\\
\bottomrule
\end{tabularx}

\begin{lstlisting}[language=tsunoo, caption={Precedência da largada (\arq{src/perfis/largada.ts})}, label={lst:largada}]
// 1. `lado_forcado` — o professor redistribuiu antes da aula. Essa decisao NAO
//    pode ser desfeita por um recalculo de afinidade no login.
// 2. Afinidade do perfil — o mecanismo normal.
// 3. Sorteio balanceado — sem perfil, ou perfil `misto`. Escolhe o lado com
//    menos runs, para a turma ainda sair perto de 50/50.
//
// Existe como modulo proprio porque DUAS rotas precisam da MESMA regra: o
// primeiro inicio e o reset da sessao. Quando o reset tinha regra propria, ele
// mantinha o lado antigo — e reaproveitar o aparelho repetia a semente sem
// ninguem perceber.
\end{lstlisting}

\begin{tabularx}{\linewidth}{@{}L{28mm}X@{}}
\toprule
\headrow \thd{Verbo} & \thd{Significado documentado}\\
\midrule
\cd{attempted} & entrou no aparelho\\
\zebra \cd{progressed} & mudou de fase\\
\cd{completed} & concluiu a missão (dossiê visto)\\
\bottomrule
\end{tabularx}

\section{exp-02 --- UnooTube}

Clone de plataforma de vídeo, para o 4º ano. O aluno salva até três vídeos; ao
terceiro, o ranking desbloqueia; \enquote{adicionar credibilidade} reordena o
ranking por fontes e é a virada pedagógica.

\begin{tabularx}{\linewidth}{@{}L{36mm}X@{}}
\toprule
\headrow \thd{Decisão} & \thd{Motivo}\\
\midrule
Gabarito (\cd{factual}) fora do \emph{bundle} & Vive em
  \arq{src/unootube/catalogo.ts}. A API corrige e devolve inscritos e fontes ---
  nunca a resposta (regra 3).\\
\zebra Ranking com colegas \textbf{fictícios} & Distribuição desenhada para a
  lição. Não há pontuação de menor real exposta a outro menor (LGPD).\\
URLs das fontes como \textbf{texto} & O \emph{sandbox} do \cd{iframe} bloqueia
  \emph{popups}, e o aluno não sai da plataforma.\\
\zebra Salvos \textbf{congelados} depois de \enquote{adicionar credibilidade} &
  O \emph{score} da experiência já foi reportado --- nota e estado não podem
  divergir. Curtir segue livre; reiniciar descongela.\\
\cd{ranking\_desbloqueado} é \emph{latch} & Remover um salvo depois não tranca o
  ranking de novo. Estado que anda para trás confunde a criança e o professor.\\
\zebra Sem \cd{viewport} & Foi desenhada para \emph{desktop}: não entra em casca
  de aparelho.\\
Tema escolhido pelo aluno & Perfil \seta{} APARÊNCIA, lembrado em
  \cd{localStorage} --- que a regra 4 proíbe para \emph{token}, não para
  preferência de tela.\\
\bottomrule
\end{tabularx}

\begin{nota}
A exp-02 é o caso canônico da \textbf{linha da revelação}
(Capítulo~\ref{cap:ds}): \emph{feed}, cromo, \emph{playlists} e canal usam a
seção do clone; ranking e os dois diálogos da missão usam a seção da plataforma.
O \emph{toast} não cruza a linha --- ele fala de dentro da ficção
(\enquote{Vídeo salvo!}).
\end{nota}

\section{exp-70 --- Vila do Lume}

Jogo 3D em primeira pessoa, \textbf{transversal} (\cd{series: []}): não pertence a
ano nenhum. É o \enquote{Ato 1} que alimenta a largada da exp-01.

\subsection{Duas leituras de uma partida, e elas não se tocam}

\begin{regra}
O retrato Schwartz (\arq{mapeamento.ts}: 10 valores + 2 eixos) lê \emph{escolhas
com custo} dentro do jogo. A missão \enquote{A Noite Longa}
(\arq{mensagens.ts}) lê \emph{julgamento sobre 5 mensagens}. A separação é
\textbf{de medida}: um golpe que custasse fagulha faria a criança que caiu
aparecer como \emph{menos generosa} no retrato, e nada no relatório mostraria
isso.
\end{regra}

\subsection{O servidor recalcula; nunca aceita o perfil do cliente}

\begin{regra}
O jogo sabe computar o retrato e \textbf{deliberadamente não o envia}:
\cd{CHAVES\_DE\_PERFIL} em \arq{src/vila/contrato.ts} \emph{recusa} \emph{payload}
que traga o perfil. É isso que permite reprocessar o histórico quando
\cd{VERSAO\_MAPEAMENTO} subir (hoje 3, herdada do piloto para os dados serem
comparáveis).
\end{regra}

\subsection{\cd{grau} é invariante de API, não regra de tela}

\begin{figure}[H]
\centering
\fitwidth{%
\begin{tikzpicture}[x=1cm, y=1cm]

  \node[classe, text width=42mm] (sc) at (0,0)
    {\etitulo{grau: 'scored'}
     \nodepart{two}
     $\geq$ 2 sinais independentes
     \nodepart{three}
     \texttt{numero: number}\\
     entra na média\\
     \textbf{é publicado com número}};

  \node[classe, text width=42mm, right=8mm of sc] (ind)
    {\etitulo{grau: 'indicative'}
     \nodepart{two}
     exatamente 1 sinal
     \nodepart{three}
     \texttt{numero: null}\\
     descrito em TEXTO\\
     fora dos eixos e das médias};

  \node[classe, text width=42mm, below=8mm of sc] (ze)
    {\etitulo{grau: 'zero'}
     \nodepart{two}
     teve a chance e não foi
     \nodepart{three}
     \texttt{numero: null}\\
     \textbf{entra na média} como zero de verdade};

  \node[classe, text width=42mm, right=8mm of ze] (no)
    {\etitulo{grau: 'none'}
     \nodepart{two}
     a oportunidade não chegou
     \nodepart{three}
     \texttt{numero: null}\\
     \textbf{sai da conta} --- não é zero};

  \node[actnote, text width=42mm, right=8mm of ind] (n1)
    {\textbf{Por que o tipo, e não a tela:} no aplicativo de origem esta decisão
     vivia na camada de interface, e qualquer tela nova que lesse os valores
     direto podia imprimir \enquote{Poder 100} a partir de um sinal só. O arquivo
     avisava disso em prosa --- e prosa não é controle.};

  \node[actnote, text width=42mm, right=8mm of no] (n2)
    {\textbf{A distinção mais importante do modelo} é entre \texttt{zero} e
     \texttt{none}. Tratar \enquote{não teve chance} como \enquote{escolheu não}
     faz uma criança parecer ter recusado algo que nunca lhe foi oferecido --- e
     esse número vai ser lido por um adulto que decide sobre ela.};

  \draw[actnoteline] (ind) -- (n1);
  \draw[actnoteline] (no) -- (n2);

\end{tikzpicture}}
\caption[Os quatro graus de medida]{Os quatro graus de medida da exp-70. A API
  devolve uma união discriminada em que \textbf{só} a variante \cd{scored} tem
  \cd{numero: number}; as outras têm \cd{null} \emph{por tipo}. O hub fica
  estruturalmente incapaz de imprimir número onde não há medida.}
\label{fig:graus}
\end{figure}

\subsection{A bolha é do servidor}

\begin{regra}
\rota{POST /vila/sessoes} atribui a bolha A/B/C \textbf{balanceada dentro da
turma}, persiste em \cd{exp\_70\_bolhas} (estável entre partidas) e devolve os
cartões --- \textbf{só os desta criança e sem gabarito}.

Antes era \cd{alunoSlot \% 3} no navegador, o que punha no \emph{bundle} a
condição do próprio instrumento e deixava as três versões legíveis antes da aula.

Nos anos iniciais (1EF--5EF) a bolha B \textbf{não é servida}
(\enquote{suporte} exige entender o que é suporte) --- regra que só a API pode
aplicar, porque só ela sabe a série.

O \emph{lock} consultivo por turma não é zelo: a sala clica junto, e sem ele os
25 caem todos na bolha A.
\end{regra}

\begin{atencao}
\cd{cartoesServidos} $\neq$ \enquote{apareceram}. O primeiro é quantos a API
entregou; o segundo, quantos chegaram a aparecer na partida. Zero aparecidos
\implica{} a tela diz \enquote{as mensagens não chegaram nesta partida},
\textbf{nunca} \enquote{ela ignorou tudo} --- é a mesma distinção entre
\cd{zero} e \cd{none}.
\end{atencao}

\subsection{Telas do professor}

\begin{tabularx}{\linewidth}{@{}L{40mm}X@{}}
\toprule
\headrow \thd{Rota} & \thd{Conteúdo e regra}\\
\midrule
\rota{.../vila-do-lume} & Ficha: para que serve, quem consome o perfil, passo a
  passo.\\
\zebra \rota{.../vila-do-lume/guia} & O debate. Projeta as mensagens das
  \textbf{três} bolhas lado a lado --- é a revelação da missão, e a primeira vez
  que a turma vê as outras versões.\\
\rota{.../resultados?turmaId=} & Listagem da turma, com o botão de zerar.\\
\zebra \rota{.../alunos/:id} & A leitura individual, que abre pelo
  \textbf{nível}.\\
\bottomrule
\end{tabularx}

\begin{regra}
\begin{enumerate}
  \item \textbf{O \rota{/guia} precisa saber a TURMA.} Nos anos iniciais a API não
    serve a bolha B, e uma tela que projetasse as três mandaria o professor
    esperar a reação de um grupo que a sala dele não tem.
    \rota{GET /vila/professor/cartoes?turmaId=} filtra pela \emph{mesma} regra de
    \cd{abrir} (\cd{bolhasElegiveis}).
  \item \textbf{Os textos vêm da API, nunca de cópia no hub.} Os textos
    \emph{são} o instrumento, e uma vírgula divergente entre o que o jogo serviu e
    o que o professor projeta põe a sala discutindo uma mensagem que ninguém
    recebeu.
  \item \textbf{A listagem da turma não mostra valor de Schwartz nenhum.}
    Listagem é ordenável, e uma coluna \enquote{Poder} viraria ranking de
    crianças. Os valores vivem na tela individual, uma criança por vez, com o
    grau ao lado.
  \item \textbf{O agregado abre para gestor; a linha de uma criança, não.} São
    \emph{controllers} separados, e não um \cd{@Roles} misto com exceção por
    método.
  \item \textbf{Zerar a turma apaga as BOLHAS junto.} Rodar de novo com a mesma
    bolha serviria à criança o mesmo golpe: ela reconhece, e a segunda rodada não
    mede nada. O que \textbf{não} é tocado é
    \cd{exp\_01\_exp\_70\_perfis\_importados} --- aquilo pertence à atividade da
    exp-01, e apagar de lá trocaria os lados de uma aula em andamento.
\end{enumerate}
\end{regra}

\subsection{A exceção de ADR: o jogo em JavaScript clássico}

\begin{decisao}
\arq{apps/experiencias/exp-70-vila-do-lume/public/jogo/} --- cerca de dez mil
linhas de JavaScript clássico com variáveis globais compartilhadas e
\arq{three.min.js} vendorizado --- \textbf{fica como está}, e está nos
\cd{ignores} do ESLint (\adrr{ADR-18}).

É um instrumento com distâncias e tempos calibrados, e o compilador não tem como
saber que uma refatoração mudou a \emph{medida}. O \emph{shim} (\arq{src/}, TS
\emph{strict}) é a única porta: \cd{window.VILA}. A rede que substitui o
\cd{tsc} ali é o \emph{smoke} visual.
\end{decisao}

\section{exp-90 --- Gerador de plano de aula}

Ferramenta do professor. Detalhada nas Seções~\ref{cap:ia} e~\ref{cap:casosdeuso};
o que resta registrar é a divisão de responsabilidade de armazenamento.

\begin{tabularx}{\linewidth}{@{}L{34mm}X@{}}
\toprule
\headrow \thd{Artefato} & \thd{Onde vive, e por quê}\\
\midrule
\arq{gerado.md} & \emph{Bucket} privado, \textbf{imutável}. É o original da IA:
  restaurar é voltar a ele.\\
\zebra \arq{atual.html} & \emph{Bucket} privado. O que o editor salvou,
  \textbf{sanitizado por \emph{allowlist}}. Versionamento de objetos dá histórico
  de cada \enquote{Salvar}.\\
Referências curriculares & PostgreSQL (\cd{exp\_90\_referencias}), como
  \textbf{texto}. Catálogo do contrato, sem \cd{escola\_id}, fora da RLS.
  Importadas por \arq{scripts/importar-referencias-plano.mjs} com
  \env{PLANOS\_IMPORT\_TOKEN}.\\
\zebra Metadados e estado & PostgreSQL (\cd{exp\_90\_planos}).\\
\bottomrule
\end{tabularx}

\begin{atencao}
\env{PLANOS\_IMPORT\_TOKEN} ausente \implica{} a rota de importação responde
\textbf{503}, não fica aberta. O mesmo vale para \env{PERFIL\_IMPORT\_TOKEN}.
\end{atencao}

\section{exp-91 --- Check-in de bem-estar}

Ferramenta do professor, detalhada na Seção~\ref{cap:ia}. O que resta registrar é
a mecânica de microdoses e a decisão de escopo de leitura.

\begin{tabularx}{\linewidth}{@{}L{34mm}X@{}}
\toprule
\headrow \thd{Mecânica} & \thd{Implementação}\\
\midrule
21 itens em 7 \emph{check-ins} & \cd{ordem\_itens} é a permutação de 1..21
  sorteada na abertura do ciclo; cada \emph{check-in} leva 3 itens, um de cada
  subescala.\\
\zebra Um \emph{check-in} por dia & \cd{UNIQUE (escola\_id, professor\_id, dia)},
  com \cd{dia} no fuso \cd{America/Sao\_Paulo}.\\
Ciclo expira em 21 dias & \cd{expira\_em}; ciclo não fechado vira
  \cd{expirado}, sem pontuação.\\
\zebra Resumo leve para o \emph{card} do hub & \rota{GET /bemestar/resumo}, sem
  IA. É o que a home do professor lê --- e a home não pode depender de uma
  chamada de modelo.\\
Apagar tudo & \rota{DELETE /bemestar/dados}.\\
\bottomrule
\end{tabularx}

\begin{decisao}
\textbf{Nenhuma rota de gestor nem agregado} (\adrr{ADR-10}). Dado de saúde
(LGPD art.~11) com RLS \emph{e} posse por \cd{professor\_id} em toda
\emph{query}; só o professor lê. Uma visão agregada e anonimizada para a rede é
decisão \emph{futura}, com DPA e \emph{k}-anonimato --- e não sai de um
\cd{GROUP BY}.
\end{decisao}

\section{Como criar uma experiência nova}

\begin{lstlisting}[language=bashunoo, caption={O caminho completo}, label={lst:novaexp}]
# 0. Uma branch por experiencia, com o nome do slug
git switch -c exp-07-fracoes

# 1. Scaffolding do _template (copia e ajusta nomes)
pnpm new:exp exp-07-fracoes

# 2. Implementar em src/App.tsx+
#    - auth, api e telemetria: SOMENTE pela SDK
#    - interface: SOMENTE o @unooria/design-system, pela ponte packages/ui,
#      na secao do clone que a experiencia imita
#    Falta componente? Ele nasce no repo storybook-unooria e chega por versao
#    nova — nunca como arquivo local.

# 3. Declarar capabilities e verbos no public/experiencia.json

# 4. Documentar verbo/objeto de telemetria no README da experiencia

# 5. Build gera dist/ 100% estatico. Testar embedado:
pnpm --filter @edu/exp-07-fracoes dev     # porta nova em lib/experiencias-dev.ts
# e abrir a bancada do hub:
#   http://localhost:3000/dev-embed?src=http://localhost:5178
\end{lstlisting}

\begin{tabularx}{\linewidth}{@{}L{34mm}X@{}}
\toprule
\headrow \thd{Checklist} & \thd{Verificação}\\
\midrule
Módulo de domínio na API & Necessário sempre que houver gabarito, estado de
  missão ou visão de professor. Prefixo minúsculo igual ao slug
  (\arq{src/<slug>/}), com \arq{catalogo.ts} + \arq{regras.ts} + \emph{service}
  com \cd{comTenancy} + \emph{controllers} por papel.\\
\zebra Tabelas & Prefixo \cd{exp\_NN\_}, \cd{escola\_id NOT NULL} + RLS.\\
Porta de \emph{dev} & Registrada em \arq{apps/hub/lib/experiencias-dev.ts}.\\
\zebra \cd{allowedHosts} do Vite & Deve incluir \cd{host.docker.internal},
  senão o hub em contêiner não lê o manifesto.\\
\cd{viewport} & Só se a experiência exigir casca de aparelho.\\
\zebra \cd{permissoes} & Só o que o \cd{iframe} realmente precisa.\\
Mídia $>$ \SI{2}{\mega\byte} & Vai para o \emph{bucket} de \emph{assets},
  referenciada por variável de \emph{build}. Nunca no \emph{git}, nunca no
  \emph{bundle}.\\
\zebra \emph{Smoke} visual & Um \emph{walkthrough} em \arq{scripts/} se a
  experiência tem tela nova no hub.\\
\bottomrule
\end{tabularx}

\begin{nota}
O \arq{\_template} continua \textbf{cru de propósito}: o que ele não traz, o
\emph{design system} traz. Um \emph{template} rico em componentes locais seria a
primeira violação da regra de interface, replicada 48 vezes.
\end{nota}
